\documentclass[10pt,a4paper]{amsart}
\usepackage[margin=19mm]{geometry}
\usepackage{amsmath,amssymb,mathtools}
\usepackage[scr=rsfs,cal=euler]{mathalfa}
\usepackage{xfrac}
\usepackage[unicode,hidelinks,bookmarks=false]{hyperref}
\newcommand{\ie}{i.e.\ }
\newcommand{\cf}{cf.\ }
\newcommand{\resp}{respectively\ }
\newcommand{\Subsection}{\subsection}
\providecommand{\nobreakdash}{\nobreak\hbox{-}}
\setlength{\emergencystretch}{2em}
\multlinegap0pt
\usepackage{float}
\usepackage{amssymb}
\usepackage{enumerate}
\usepackage{xcolor}
\newtheorem{proposition}{Proposition}
\newtheorem{corollary}{Corollary}
\newtheorem{theorem}{Theorem}
\title{On Some Applications to the Newton--Raphson Method}
\author{Treanungkur Mal}
\date{Source: arXiv:2211.07433v2 (CC BY 4.0)}
\begin{document}
\maketitle

\noindent\emph{Source and licence.} On Some Applications to the Newton--Raphson Method. Version arXiv:2211.07433v2 (CC BY 4.0), distributed by arXiv under the Creative Commons Attribution 4.0 International licence, http://creativecommons.org/licenses/by/4.0/, as recorded in the arXiv licence metadata for this exact version and shown on the abstract page https://arxiv.org/abs/2211.07433v2. Full licence notice: \texttt{licenses/arxiv-2211.07433-CC-BY-4.0.txt}.\par\smallskip

\noindent\textit{Editorial scope.} Counted unit: the article's corollary corollary (source0.tex lines 485--505). The article's complete original proof of it is delivered with it (source0.tex lines 507--628, one \texttt{proof} environment of 122 lines). No mathematics is written, repaired or re-derived: the statement and the proof are the article's own lines, in the article's own notation, and no retained line is substituted or re-flowed.  Retained uncounted as the source-local context the proof needs: lines 205--219; lines 279--298; lines 356--367; lines 406--429.  Nothing was omitted from any retained span, so the package carries no omission record.  No retained line contains a citation, so the wrapper carries no bibliography and no bibliography entry.  No retained line contains a figure environment, an image-inclusion command or drawing code, so no image is delivered and the package declares no asset dependency.  Editorial formatting only, in addition: an AMS \texttt{amsart} preamble that restates the article's own declarations and macros used by the retained lines, namely the environments \texttt{proposition}, \texttt{corollary}, \texttt{theorem} on separate counters, together with the standard packages those lines need.  The retained environments print in this document as Proposition 1, Corollary 1, Theorem 1 in order of appearance, whereas the article prints the same items with the numbering of its own full document.  The complete native source of the article as distributed in the arXiv e-print for this exact version is bundled unchanged as \texttt{sources/133456/source0.tex}.
\begin{proposition}\label{prop: 2,1}
For every $n\ge0$,
\[
e_{n+1}
=\frac{f''(\xi_n)}{2f'(x_n)}\,e_n^2,
\]
where $\xi_n\in(a,x_n)$. Since $f'$ is continuous and positive on the compact interval $\Omega$, there exists $c:=\min_{x\in\Omega}f'(x)>0$. Consequently,
\[
e_{n+1}\le\frac{M}{2c}e_n^2.
\]
Moreover,
\[
h_{n+1}\le\frac{M}{2c}h_n^2.
\]
\end{proposition}

\begin{proposition}[Exact quadrature decomposition]\label{prop: 3.1}
For every $N\ge1$,
\[
\int_a^bf(x)\,dx
=
Q_N
-
\sum_{n=0}^{N-1}E_n
+
R_N,
\]
where
\[
E_n
=
T_n-\int_{x_{n+1}}^{x_n}f(x)\,dx,
\qquad
R_N=\int_a^{x_N}f(x)\,dx.
\]
\end{proposition}

\begin{corollary}\label{cor: 3.3}
For every Newton interval,
\[
\frac{m}{12}h_n^3
\le
E_n
\le
\frac{M}{12}h_n^3.
\]
In particular, each local trapezoidal approximation overestimates a strictly
convex function.
\end{corollary}

\begin{theorem}\label{thm: 4.1}
Assume
\[
\sum_{n=0}^{\infty} h_n^3 < \infty.
\]
Then
\[
\lim_{N\to\infty} Q_N
=
\int_a^b f(x)\,dx
+
\sum_{n=0}^{\infty}E_n.
\]
In particular,
\[
\lim_{N\to\infty}
\left(
Q_N-\int_a^b f(x)\,dx
\right)
=
\sum_{n=0}^{\infty}E_n
>0.
\]
\end{theorem}

\begin{corollary}
Under the standing assumptions on $f$, the sequence $(Q_N)$ converges and
\[
\lim_{N\to\infty} Q_N
=
\int_a^b f(x)\,dx + E_\infty,
\qquad
E_\infty:=\sum_{n=0}^{\infty}E_n.
\]
Moreover, there exist constants $C>0$ and $N_0\in\mathbb{N}$ such that, for
all $N\geq N_0$,
\[
\left|
Q_N-
\left(
\int_a^b f(x)\,dx+E_\infty
\right)
\right|
\leq C e_N^2.
\]
\end{corollary}

\begin{proof}
By Theorem~\ref{thm: 4.1},
the series $\sum_{n=0}^{\infty}E_n$ converges and
\[
\lim_{N\to\infty}Q_N
=
\int_a^b f(x)\,dx+E_\infty.
\]
From the exact decomposition in Proposition~\ref{prop: 3.1},
\[
Q_N-
\left(
\int_a^b f(x)\,dx+E_\infty
\right)
=
R_N-\sum_{n=N}^{\infty}E_n.
\]

Let
\[
L:=\max_{x\in[a,b]}|f'(x)|.
\]
Since $f(a)=0$, the mean value theorem yields
\[
|f(x)|
=
|f(x)-f(a)|
\leq L(x-a),
\qquad x\in[a,b].
\]
Hence
\[
|R_N|
=
\left|\int_a^{x_N}f(x)\,dx\right|
\leq
\int_a^{x_N}L(x-a)\,dx
=
\frac{L}{2}e_N^2.
\]

On the other hand, Proposition~\ref{prop: 2,1} gives
\[
h_{n+1}\leq \frac{M}{2c}h_n^2.
\]
Since $h_n\to0$, there exists $N_0\in\mathbb{N}$ such that
\[
\frac{M}{2c}h_n\leq\frac12,
\qquad n\geq N_0.
\]
Thus
\[
h_{n+1}\leq\frac12h_n,
\qquad n\geq N_0.
\]
It follows that, for $N\geq N_0$,
\[
\sum_{n=N}^{\infty}h_n^3
\leq
h_N^3\sum_{k=0}^{\infty}\frac{1}{2^{3k}}
=
\frac{8}{7}h_N^3.
\]
By Corollary~\ref{cor: 3.3},
\[
\sum_{n=N}^{\infty}E_n
\leq
\frac{M}{12}\sum_{n=N}^{\infty}h_n^3
\leq
\frac{2M}{21}h_N^3.
\]
Since
\[
h_N=x_N-x_{N+1}\leq x_N-a=e_N,
\]
we obtain
\[
\sum_{n=N}^{\infty}E_n
\leq
\frac{2M}{21}e_N^3.
\]
Therefore, for $N\geq N_0$,
\[
\begin{aligned}
\left|
Q_N-
\left(
\int_a^b f(x)\,dx+E_\infty
\right)
\right|
&\leq
|R_N|+\sum_{n=N}^{\infty}E_n\\
&\leq
\frac{L}{2}e_N^2
+
\frac{2M}{21}e_N^3.
\end{aligned}
\]
Since $e_N\leq b-a$,
\[
\left|
Q_N-
\left(
\int_a^b f(x)\,dx+E_\infty
\right)
\right|
\leq
\left(
\frac{L}{2}
+
\frac{2M}{21}(b-a)
\right)e_N^2.
\]
Thus the asserted estimate holds with
\[
C=
\frac{L}{2}
+
\frac{2M}{21}(b-a).
\]
This completes the proof.
\end{proof}

\end{document}
